% ══════════════════════════════════════════════════════════════
%  Rapport PFE — Météo Tunisie · Pipeline Big Data
%  ISIMS Sfax · Département ADBD · 2025-2026
% ══════════════════════════════════════════════════════════════
\documentclass[a4paper,12pt,oneside]{report}

% ─── Encodage et langue ───
\usepackage[utf8]{inputenc}
\usepackage[T1]{fontenc}
\usepackage[french]{babel}
\usepackage{csquotes}

% ─── Mise en page ───
\usepackage[top=2.5cm,bottom=2.5cm,left=3cm,right=2.5cm]{geometry}
\usepackage{setspace}
\onehalfspacing
\usepackage{fancyhdr}
\pagestyle{fancy}
\fancyhf{}
\fancyhead[L]{\leftmark}
\fancyhead[R]{\thepage}
\renewcommand{\headrulewidth}{0.4pt}

% ─── Graphiques et couleurs ───
\usepackage{graphicx}
\usepackage[dvipsnames,table]{xcolor}
\usepackage{float}
\usepackage{subcaption}
\usepackage{tikz}
\usetikzlibrary{shapes,arrows.meta,positioning,fit}

% ─── Tableaux ───
\usepackage{booktabs}
\usepackage{tabularx}
\usepackage{longtable}
\usepackage{multirow}
\usepackage{array}

% ─── Code source ───
\usepackage{listings}
\lstset{
  basicstyle=\ttfamily\small,
  backgroundcolor=\color{gray!8},
  frame=single,
  rulecolor=\color{gray!40},
  breaklines=true,
  numbers=left,
  numberstyle=\tiny\color{gray},
  keywordstyle=\color{blue!70!black}\bfseries,
  commentstyle=\color{green!50!black}\itshape,
  stringstyle=\color{red!60!black},
  showstringspaces=false,
  tabsize=2,
  captionpos=b
}

% ─── Liens et PDF ───
\usepackage[hidelinks]{hyperref}
\hypersetup{
  pdftitle={Météo Tunisie — Pipeline Big Data},
  pdfauthor={Équipe Projet ISIMS},
  pdfsubject={Rapport PFE},
}

% ─── Divers ───
\usepackage{amsmath,amssymb}
\usepackage{enumitem}
\usepackage{pifont}
\usepackage{tcolorbox}
\tcbuselibrary{skins,breakable}
\usepackage{caption}
\usepackage{url}

% ─── Commandes personnalisées ───
\newcommand{\imgplaceholder}[2]{%
  \begin{tcolorbox}[colback=blue!5,colframe=blue!40,title=#1]
    \centering\itshape #2\\[6pt]
    \texttt{$\rightarrow$ Insérer la capture ici}
  \end{tcolorbox}
}

% ══════════════════════════════════════════════════════════════
\begin{document}

% ─────────────────────────────────────────
%  PAGE DE GARDE
% ─────────────────────────────────────────
\begin{titlepage}
\centering
\vspace*{1cm}

% Logo université
% \includegraphics[width=4cm]{figures/logo_isims.png}
\imgplaceholder{Logo}{Logo ISIMS Sfax}

\vspace{0.5cm}
{\large \textbf{Université de Sfax}}\\[4pt]
{\large Institut Supérieur d'Informatique et de Multimédia de Sfax}\\[4pt]
{\large Département ADBD}\\[1.5cm]

{\Large \textbf{RAPPORT DE PROJET}}\\[0.3cm]
{\large Année universitaire 2025--2026}\\[1.5cm]

\rule{\textwidth}{1.5pt}\\[0.4cm]
{\LARGE \textbf{Météo Tunisie}}\\[0.2cm]
{\Large Pipeline Big Data de Collecte, Analyse et\\Visualisation des Données Météorologiques}\\[0.4cm]
\rule{\textwidth}{1.5pt}\\[2cm]

\begin{minipage}[t]{0.45\textwidth}
\textbf{Réalisé par :}\\
Mohamed Chaari\\
Yessine [Nom]\\
Moataz [Nom]\\
Brahim [Nom]
\end{minipage}
\hfill
\begin{minipage}[t]{0.45\textwidth}
\raggedleft
\textbf{Encadré par :}\\
Mme Nesrine [Nom]
\end{minipage}

\vfill
{\large Sfax, 2025--2026}
\end{titlepage}

% ─────────────────────────────────────────
%  DÉDICACES
% ─────────────────────────────────────────
\chapter*{Dédicaces}
\addcontentsline{toc}{chapter}{Dédicaces}
\thispagestyle{empty}
\vspace*{3cm}
\begin{center}
\textit{À nos familles,\\
qui nous ont soutenus tout au long de ce parcours.\\[1cm]
À nos enseignants,\\
pour leur patience et leur dévouement.\\[1cm]
À tous ceux qui ont contribué,\\
de près ou de loin, à la réalisation de ce projet.}
\end{center}
\newpage

% ─────────────────────────────────────────
%  REMERCIEMENTS
% ─────────────────────────────────────────
\chapter*{Remerciements}
\addcontentsline{toc}{chapter}{Remerciements}

Nous tenons à exprimer nos sincères remerciements à notre encadrante, \textbf{Mme Nesrine [Nom]}, pour son accompagnement, ses conseils précieux et sa disponibilité tout au long de ce projet.

Nous remercions également l'ensemble du corps enseignant du département ADBD de l'ISIMS de Sfax pour la formation de qualité qu'ils nous ont dispensée.

Enfin, nous adressons nos remerciements à toutes les personnes qui ont contribué, de près ou de loin, à la réalisation de ce travail.

\newpage

% ─────────────────────────────────────────
%  RÉSUMÉ
% ─────────────────────────────────────────
\chapter*{Résumé}
\addcontentsline{toc}{chapter}{Résumé}

\textbf{Français :}\\
Ce projet porte sur la conception et le développement d'une plateforme Big Data dédiée à la collecte, l'analyse et la visualisation des données météorologiques de 221 villes tunisiennes. L'architecture s'appuie sur un pipeline de streaming en temps réel utilisant Apache Kafka (Aiven) pour l'ingestion, Supabase (PostgreSQL) pour le stockage, Apache Airflow pour l'orchestration des analyses, et une API REST FastAPI pour l'exposition des données. Le système couvre l'historique météorologique de 2020 à 2024, l'ingestion temps réel toutes les 15 minutes, ainsi que des analyses avancées incluant la détection de pics thermiques, les corrélations saisonnières et les tendances annuelles.

\vspace{0.8cm}
\textbf{English:}\\
This project focuses on designing and developing a Big Data platform for collecting, analyzing, and visualizing meteorological data from 221 Tunisian cities. The architecture relies on a real-time streaming pipeline using Apache Kafka (Aiven) for ingestion, Supabase (PostgreSQL) for storage, Apache Airflow for analysis orchestration, and a FastAPI REST API for data exposure. The system covers the weather history from 2020 to 2024, real-time ingestion every 15 minutes, and advanced analytics including thermal peak detection, seasonal correlations, and annual trends.

\newpage

% ─────────────────────────────────────────
%  LISTE DES ABRÉVIATIONS
% ─────────────────────────────────────────
\chapter*{Liste des Abréviations}
\addcontentsline{toc}{chapter}{Liste des Abréviations}

\begin{tabularx}{\textwidth}{lX}
\toprule
\textbf{Abréviation} & \textbf{Signification} \\
\midrule
API     & Application Programming Interface \\
ADBD    & Administration des Bases de Données \\
DAG     & Directed Acyclic Graph \\
DB      & Database (Base de données) \\
E2E     & End-to-End \\
JWT     & JSON Web Token \\
KPI     & Key Performance Indicator \\
OWM     & OpenWeatherMap \\
PaaS    & Platform as a Service \\
REST    & Representational State Transfer \\
SaaS    & Software as a Service \\
SQL     & Structured Query Language \\
SSL     & Secure Sockets Layer \\
UI/UX   & User Interface / User Experience \\
US      & User Story \\
\bottomrule
\end{tabularx}

\newpage

% ─────────────────────────────────────────
%  TABLES
% ─────────────────────────────────────────
\tableofcontents
\newpage
\listoffigures
\addcontentsline{toc}{chapter}{Liste des Figures}
\newpage
\listoftables
\addcontentsline{toc}{chapter}{Liste des Tableaux}
\newpage

% ══════════════════════════════════════════════════════════════
%  CHAPITRES
% ══════════════════════════════════════════════════════════════
% ══════════════════════════════════════════
%  CHAPITRE 1 — INTRODUCTION GÉNÉRALE
% ══════════════════════════════════════════
\chapter{Introduction Générale}

\section{Contexte et motivation}
La Tunisie, pays méditerranéen aux contrastes climatiques marqués, est exposée à des phénomènes météorologiques variés~: canicules estivales, précipitations intenses et vents violents. La compréhension fine de ces phénomènes est essentielle pour l'agriculture, le tourisme et la gestion des risques naturels. Or, les données météorologiques existantes sont souvent dispersées, difficilement accessibles et insuffisamment exploitées.

Dans ce contexte, les technologies Big Data offrent une opportunité unique de centraliser, analyser et valoriser ces données à grande échelle. L'émergence des services Cloud managés (SaaS/PaaS) tels que Supabase, Aiven et Astronomer permet de déployer des architectures de données sophistiquées sans infrastructure on-premise.

\section{Problématique}
Comment concevoir et déployer un pipeline de données météorologiques de bout en bout, capable de~:
\begin{itemize}
  \item Collecter les données historiques et temps réel de \textbf{221 villes tunisiennes}~?
  \item Ingérer ces données de manière fiable via un système de streaming distribué~?
  \item Orchestrer des analyses statistiques avancées de manière automatisée~?
  \item Exposer les résultats à travers une interface web interactive et intuitive~?
\end{itemize}

\section{Objectifs du projet}
Les objectifs principaux de ce projet sont~:
\begin{enumerate}
  \item \textbf{Infrastructure Cloud}~: Mettre en place une base PostgreSQL sur Supabase et un broker Kafka sur Aiven pour l'ingestion en temps réel.
  \item \textbf{Pipeline de données}~: Développer des Producers et Consumers Python pour alimenter la base de données à partir des APIs météorologiques (Open-Meteo, OpenWeatherMap).
  \item \textbf{Orchestration analytique}~: Créer des DAGs Apache Airflow pour automatiser les analyses mensuelles, annuelles, la détection de pics et les corrélations.
  \item \textbf{API et Visualisation}~: Exposer les données via une API REST (FastAPI) et un dashboard interactif avec cartes et graphiques.
  \item \textbf{Déploiement Cloud}~: Déployer l'ensemble sur des services managés (Astronomer pour Airflow, Render pour l'API, GitHub Pages pour le frontend).
\end{enumerate}

\section{Structure du rapport}
Ce rapport est organisé comme suit~:
\begin{itemize}
  \item \textbf{Chapitre 2}~: État de l'art sur les systèmes météorologiques et les architectures Big Data.
  \item \textbf{Chapitre 3}~: Méthodologie Scrum et spécifications globales du projet.
  \item \textbf{Chapitre 4}~: Sprint 1 — Infrastructure Cloud et flux temps réel.
  \item \textbf{Chapitre 5}~: Sprint 2 — Analyses avancées, alertes et déploiement.
  \item \textbf{Chapitre 6}~: Tests globaux et validation finale.
\end{itemize}

% ══════════════════════════════════════════
%  CHAPITRE 2 — ÉTAT DE L'ART
% ══════════════════════════════════════════
\chapter{État de l'Art}

\section{Systèmes météorologiques existants}
Plusieurs plateformes offrent des données météorologiques à l'échelle mondiale~:

\begin{table}[H]
\centering
\caption{Comparaison des sources de données météorologiques}
\label{tab:sources_meteo}
\begin{tabularx}{\textwidth}{lXccc}
\toprule
\textbf{Source} & \textbf{Description} & \textbf{Gratuit} & \textbf{Historique} & \textbf{Temps réel} \\
\midrule
Open-Meteo       & API open-source, données ECMWF   & \ding{51} & \ding{51} & \ding{51} \\
OpenWeatherMap    & API freemium, large couverture   & Partiel   & \ding{51} & \ding{51} \\
Météo-France      & Service national français        & \ding{55} & \ding{51} & \ding{51} \\
Copernicus (ERA5) & Données satellite réanalyse      & \ding{51} & \ding{51} & \ding{55} \\
INM Tunisie       & Institut National de la Météo    & \ding{55} & Limité    & Limité    \\
\bottomrule
\end{tabularx}
\end{table}

Notre choix s'est porté sur \textbf{Open-Meteo} comme source principale (gratuit, sans clé API, haute résolution) et \textbf{OpenWeatherMap} comme source de secours pour les données temps réel.

\section{Technologies Big Data pour les données météo}

\subsection{Apache Kafka}
Apache Kafka est une plateforme de streaming distribué conçue pour gérer des flux de données à haut débit. Dans notre contexte, Kafka assure le découplage entre les producteurs de données météo et les consommateurs qui les persistent en base.

\subsection{Apache Airflow}
Apache Airflow est un orchestrateur de workflows qui permet de définir, planifier et surveiller des pipelines de données sous forme de DAGs (Directed Acyclic Graphs). Il est particulièrement adapté aux traitements batch analytiques.

\subsection{PostgreSQL et Supabase}
PostgreSQL est le SGBD relationnel open-source le plus avancé. Supabase fournit une couche managée (PaaS) avec API REST automatique, authentification et stockage objet.

\section{Architectures streaming (Lambda, Kappa)}

\begin{table}[H]
\centering
\caption{Comparaison des architectures de traitement de données}
\label{tab:archi}
\begin{tabularx}{\textwidth}{lXX}
\toprule
\textbf{Architecture} & \textbf{Avantages} & \textbf{Inconvénients} \\
\midrule
Lambda  & Double couche batch + stream, résultats fiables & Complexité de maintenance, code dupliqué \\
Kappa   & Architecture simplifiée, un seul pipeline stream & Reprocessing coûteux, moins adapté aux analyses lourdes \\
Hybride & Flexibilité, adapté aux cas mixtes & Nécessite une bonne orchestration \\
\bottomrule
\end{tabularx}
\end{table}

\section{Comparaison des solutions existantes}

\begin{table}[H]
\centering
\caption{Positionnement de notre solution}
\label{tab:positionnement}
\begin{tabularx}{\textwidth}{lcccc}
\toprule
\textbf{Critère} & \textbf{Notre solution} & \textbf{Windy} & \textbf{Infoclimat} & \textbf{Weather.gov} \\
\midrule
Focus Tunisie         & \ding{51} & \ding{55} & \ding{55} & \ding{55} \\
221 villes            & \ding{51} & \ding{55} & \ding{55} & \ding{55} \\
Analyses statistiques & \ding{51} & \ding{55} & Partiel   & \ding{55} \\
Détection de pics     & \ding{51} & \ding{55} & \ding{55} & \ding{55} \\
Open-source           & \ding{51} & \ding{55} & \ding{55} & \ding{51} \\
Pipeline automatisé   & \ding{51} & N/A       & \ding{55} & Partiel   \\
\bottomrule
\end{tabularx}
\end{table}

\section{Positionnement de notre solution}
Notre plateforme se distingue par sa spécialisation sur le territoire tunisien avec une couverture exhaustive de 221 villes, et par l'intégration d'analyses statistiques avancées (corrélations de Pearson, détection de pics par z-score, tendances de réchauffement) directement dans le pipeline de données orchestré par Airflow.

% ══════════════════════════════════════════
%  CHAPITRE 3 — MÉTHODOLOGIE SCRUM
% ══════════════════════════════════════════
\chapter{Méthodologie Scrum et Spécifications Globales}

\section{Présentation de la démarche Scrum}
Le projet a été mené selon la méthodologie agile \textbf{Scrum}, structurée en deux sprints de durée égale. Cette approche itérative nous a permis de livrer un incrément fonctionnel à chaque fin de sprint et d'adapter les priorités en fonction des retours.

\begin{figure}[H]
\centering
\imgplaceholder{Figure~\ref{fig:scrum_cycle} — Cycle Scrum}{Diagramme illustrant le cycle Scrum avec Product Backlog, Sprint Planning, Daily Scrum, Sprint Review et Retrospective}
\caption{Cycle de développement Scrum appliqué au projet}
\label{fig:scrum_cycle}
\end{figure}

\textbf{Rôles~:}
\begin{itemize}
  \item \textbf{Product Owner}~: Mme Nesrine [Nom] (encadrante)
  \item \textbf{Scrum Master}~: Mohamed Chaari
  \item \textbf{Équipe de développement}~: Mohamed, Yessine, Moataz, Brahim
\end{itemize}

\section{Identification des acteurs}

\begin{table}[H]
\centering
\caption{Acteurs du système}
\label{tab:acteurs}
\begin{tabularx}{\textwidth}{lX}
\toprule
\textbf{Acteur} & \textbf{Description} \\
\midrule
Utilisateur    & Consulte le dashboard, les prévisions et les analyses \\
Administrateur & Gère l'infrastructure Cloud et les configurations \\
Analyste       & Consulte les analyses statistiques avancées et les alertes \\
Système        & Exécute automatiquement les pipelines et les DAGs \\
\bottomrule
\end{tabularx}
\end{table}

\section{Besoins fonctionnels et non fonctionnels globaux}

\subsection{Besoins fonctionnels}
\begin{enumerate}
  \item Collecter les données météo historiques (2020--2024) et temps réel pour 221 villes.
  \item Ingérer les données via un système de streaming fiable (Kafka).
  \item Stocker les données dans une base PostgreSQL Cloud (Supabase).
  \item Orchestrer les analyses automatisées (moyennes, corrélations, pics, tendances).
  \item Exposer les données via une API REST documentée.
  \item Afficher les résultats dans un dashboard interactif avec cartes et graphiques.
\end{enumerate}

\subsection{Besoins non fonctionnels}
\begin{itemize}
  \item \textbf{Performance}~: Ingestion de $>$200 enregistrements/batch sans perte.
  \item \textbf{Disponibilité}~: Infrastructure Cloud avec haute disponibilité (Supabase, Aiven).
  \item \textbf{Sécurité}~: Connexions SSL pour Kafka, mots de passe chiffrés.
  \item \textbf{Maintenabilité}~: Code modulaire avec tests unitaires ($>$100 tests).
  \item \textbf{Scalabilité}~: Architecture découplée permettant l'ajout de nouvelles villes ou sources.
\end{itemize}

\section{Diagramme de cas d'utilisation global}

\begin{figure}[H]
\centering
\imgplaceholder{Figure~\ref{fig:uc_global} — Diagramme de cas d'utilisation global}{Diagramme UML montrant les 4 acteurs (Utilisateur, Administrateur, Analyste, Système) et les cas d'utilisation principaux~: Consulter Dashboard, Voir Prévisions, Analyser Tendances, Recevoir Alertes, Configurer Infrastructure, Exécuter Pipeline}
\caption{Diagramme de cas d'utilisation global}
\label{fig:uc_global}
\end{figure}

\section{Product Backlog complet}

\begin{table}[H]
\centering
\caption{Product Backlog — Liste des User Stories}
\label{tab:backlog}
\begin{tabularx}{\textwidth}{clXcc}
\toprule
\textbf{US} & \textbf{Sprint} & \textbf{Description} & \textbf{Priorité} & \textbf{Points} \\
\midrule
US1 & S1 & Configuration de Supabase et Aiven            & Haute & 8 \\
US2 & S1 & Producer live et Consumer vers la DB            & Haute & 13 \\
US3 & S1 & Endpoints API /summary et /forecast             & Moyenne & 8 \\
US4 & S1 & Authentification JWT                             & Moyenne & 5 \\
US5 & S1 & Dashboard et graphiques de prévisions            & Haute & 13 \\
US6 & S1 & Validation des flux et plan de test              & Basse & 3 \\
\midrule
US7 & S2 & DAGs Airflow (moyennes, corrélations)            & Haute & 13 \\
US8 & S2 & Détection de pics (canicule, tempêtes)            & Haute & 8 \\
US9 & S2 & Documentation, monitoring et livraison finale     & Moyenne & 5 \\
\bottomrule
\end{tabularx}
\end{table}

\section{Architecture globale du système}

\begin{figure}[H]
\centering
\imgplaceholder{Figure~\ref{fig:archi_globale} — Architecture globale du pipeline}{Diagramme d'architecture montrant~: APIs Météo $\rightarrow$ Producer Python $\rightarrow$ Aiven Kafka $\rightarrow$ Consumer Python $\rightarrow$ Supabase PostgreSQL $\rightarrow$ Airflow DAGs $\rightarrow$ FastAPI $\rightarrow$ Dashboard Frontend}
\caption{Architecture globale du système Météo Tunisie}
\label{fig:archi_globale}
\end{figure}

\section{Planification des Sprints}

\begin{table}[H]
\centering
\caption{Répartition des Sprints}
\label{tab:sprints}
\begin{tabularx}{\textwidth}{lXl}
\toprule
\textbf{Sprint} & \textbf{Objectif} & \textbf{User Stories} \\
\midrule
Sprint 1 & Infrastructure Cloud et Flux Temps Réel    & US1 à US6 \\
Sprint 2 & Analyses Avancées, Alertes et Déploiement   & US7 à US9 \\
\bottomrule
\end{tabularx}
\end{table}

% ══════════════════════════════════════════
%  CHAPITRE 4 — SPRINT 1
% ══════════════════════════════════════════
\chapter{Sprint 1~: Infrastructure Cloud et Flux Temps Réel}

\section{Introduction}
Le premier sprint marque le lancement technique du projet. L'objectif est de mettre en place une architecture robuste et évolutive en s'appuyant sur des services Cloud (SaaS/PaaS). Cette phase se concentre sur la création de la base de données sur Supabase, la configuration du broker de messages Kafka via Aiven, et le développement du socle de l'API permettant de faire circuler les premières données météorologiques de bout en bout.

\section{Objectif du Sprint et Sprint Backlog}

\subsection{Objectif attendu}
Déployer l'infrastructure Cloud, charger l'historique météo (2020--2024) pour les 221 villes, et activer le flux d'ingestion en temps réel pour alimenter un tableau de bord dynamique et sécurisé.

\subsection{Sprint Backlog}

\begin{table}[H]
\centering
\caption{Sprint Backlog — Sprint 1}
\label{tab:sprint1_backlog}
\begin{tabularx}{\textwidth}{clX}
\toprule
\textbf{US} & \textbf{Catégorie} & \textbf{Description} \\
\midrule
US1 & Infrastructure & Configuration de Supabase (PostgreSQL) et Aiven (Kafka) \\
US2 & Ingestion      & Développement du Producer live et du Consumer vers la DB \\
US3 & API Core       & Création des endpoints \texttt{/summary} et \texttt{/forecast} \\
US4 & Sécurité       & Mise en place de l'authentification JWT \\
US5 & UI/UX          & Intégration du Dashboard et des graphiques de prévisions \\
US6 & Qualité        & Validation des flux et rédaction du plan de test \\
\bottomrule
\end{tabularx}
\end{table}

\section{Analyse et Spécification}

\subsection{Diagramme de cas d'utilisation du Sprint 1}

\begin{figure}[H]
\centering
\imgplaceholder{Figure~\ref{fig:uc_sprint1} — Cas d'utilisation Sprint 1}{Diagramme UML avec acteurs Administrateur, Système et Utilisateur. Cas~: Configurer Infrastructure, Ingérer Données Historiques, Consulter Dashboard, S'authentifier}
\caption{Diagramme de cas d'utilisation — Sprint 1}
\label{fig:uc_sprint1}
\end{figure}

\subsection{Description textuelle des cas d'utilisation}

\begin{table}[H]
\centering
\caption{Cas d'utilisation — Consulter les données temps réel (US5)}
\label{tab:uc_dashboard}
\begin{tabularx}{\textwidth}{lX}
\toprule
\textbf{Élément} & \textbf{Description} \\
\midrule
Acteur         & Utilisateur \\
Pré-conditions & L'utilisateur est authentifié ; les flux Aiven/Supabase sont actifs \\
Scénario nominal & 1. L'utilisateur accède au dashboard \\
               & 2. Le système appelle l'API \texttt{/api/dashboard/summary} \\
               & 3. L'interface affiche les KPIs météo mis à jour \\
Post-conditions & Les données sont affichées avec la dernière mise à jour \\
\bottomrule
\end{tabularx}
\end{table}

\begin{table}[H]
\centering
\caption{Cas d'utilisation — S'authentifier (US4)}
\label{tab:uc_auth}
\begin{tabularx}{\textwidth}{lX}
\toprule
\textbf{Élément} & \textbf{Description} \\
\midrule
Acteur         & Utilisateur \\
Scénario nominal & 1. Saisie des identifiants (email + mot de passe) \\
               & 2. Vérification dans la table \texttt{users} de Supabase \\
               & 3. Génération et renvoi d'un token JWT \\
Scénario alternatif & Identifiants incorrects $\rightarrow$ message d'erreur \\
\bottomrule
\end{tabularx}
\end{table}

\section{Conception}

\subsection{Conception de la base de données}

\begin{figure}[H]
\centering
\imgplaceholder{Figure~\ref{fig:schema_db} — Schéma relationnel Supabase}{Schéma montrant les 8 tables~: weather\_historical, weather\_current, weather\_forecast, monthly\_averages, correlations, temperature\_peaks, annual\_stats, pipeline\_runs avec leurs relations et clés}
\caption{Schéma relationnel de la base de données Supabase}
\label{fig:schema_db}
\end{figure}

\subsection{Dictionnaire de données}

\begin{table}[H]
\centering
\caption{Dictionnaire de données — Table \texttt{weather\_historical}}
\label{tab:dict_data}
\begin{tabularx}{\textwidth}{llcX}
\toprule
\textbf{Colonne} & \textbf{Type} & \textbf{Nullable} & \textbf{Description} \\
\midrule
id            & SERIAL     & Non & Identifiant auto-incrémenté \\
date          & DATE       & Non & Date de la mesure \\
city          & VARCHAR    & Non & Nom de la ville \\
governorate   & VARCHAR    & Oui & Gouvernorat \\
region        & VARCHAR    & Oui & Région (Nord, Centre, Sud) \\
temperature   & REAL       & Oui & Température en °C \\
humidity      & REAL       & Oui & Humidité relative en \% \\
precipitation & REAL       & Oui & Précipitations en mm \\
wind\_speed   & REAL       & Oui & Vitesse du vent en km/h \\
ingested\_at  & TIMESTAMPTZ & Oui & Horodatage d'ingestion \\
\bottomrule
\end{tabularx}
\end{table}

\subsection{Diagrammes de séquence}

\begin{figure}[H]
\centering
\imgplaceholder{Figure~\ref{fig:seq_ingestion} — Séquence d'ingestion}{Diagramme de séquence~: API Open-Meteo $\rightarrow$ Producer Python $\rightarrow$ Topic Kafka (Aiven) $\rightarrow$ Consumer Python $\rightarrow$ Supabase PostgreSQL avec batch upsert}
\caption{Diagramme de séquence — Flux d'ingestion historique}
\label{fig:seq_ingestion}
\end{figure}

\begin{figure}[H]
\centering
\imgplaceholder{Figure~\ref{fig:seq_realtime} — Séquence temps réel}{Diagramme de séquence~: Scheduler (15 min) $\rightarrow$ Producer $\rightarrow$ Kafka topic weather-current $\rightarrow$ Consumer $\rightarrow$ DB $\rightarrow$ API $\rightarrow$ Dashboard auto-refresh}
\caption{Diagramme de séquence — Flux temps réel (rafraîchissement 15 min)}
\label{fig:seq_realtime}
\end{figure}

\subsection{Prototypes d'interfaces}

\begin{figure}[H]
\centering
\imgplaceholder{Figure~\ref{fig:wireframe_login} — Wireframe Login}{Maquette de la page login.html~: formulaire centré avec champs email et mot de passe, bouton de connexion, thème sombre}
\caption{Wireframe — Page de connexion}
\label{fig:wireframe_login}
\end{figure}

\begin{figure}[H]
\centering
\imgplaceholder{Figure~\ref{fig:wireframe_dashboard} — Wireframe Dashboard}{Maquette du dashboard~: grille de cartes KPI (température, humidité, vent), carte Leaflet.js, graphiques temporels Plotly.js}
\caption{Wireframe — Dashboard principal}
\label{fig:wireframe_dashboard}
\end{figure}

\section{Implémentation}

\subsection{Environnement de développement}

\begin{table}[H]
\centering
\caption{Environnement technique — Sprint 1}
\label{tab:env_sprint1}
\begin{tabularx}{\textwidth}{lX}
\toprule
\textbf{Composant} & \textbf{Technologie} \\
\midrule
Langage          & Python 3.10 \\
Base de données  & PostgreSQL 15 via Supabase (Session Pooler) \\
Broker messages  & Apache Kafka 3.x via Aiven (SSL) \\
API              & FastAPI + Uvicorn \\
Frontend         & HTML5, CSS3, JavaScript, Plotly.js, Leaflet.js \\
Bibliothèques   & kafka-python-ng, psycopg2, SQLAlchemy, pandas \\
Versioning       & Git + GitHub \\
\bottomrule
\end{tabularx}
\end{table}

\subsection{Infrastructure Kafka sur Aiven}

\begin{figure}[H]
\centering
\imgplaceholder{Figure~\ref{fig:aiven_console} — Console Aiven}{Capture d'écran de la console Aiven montrant le service Kafka avec les topics créés et la configuration SSL}
\caption{Console Aiven — Service Kafka configuré}
\label{fig:aiven_console}
\end{figure}

Trois topics ont été créés~:
\begin{itemize}
  \item \texttt{weather-historical}~: données historiques (2020--2024)
  \item \texttt{weather-current}~: mesures temps réel (toutes les 15 min)
  \item \texttt{weather-forecast}~: prévisions à 7 jours
\end{itemize}

\subsection{Base de données Supabase}

\begin{figure}[H]
\centering
\imgplaceholder{Figure~\ref{fig:supabase_tables} — Tables Supabase}{Capture d'écran de l'éditeur SQL Supabase montrant les tables créées et le nombre de lignes}
\caption{Tables de la base de données sur Supabase}
\label{fig:supabase_tables}
\end{figure}

\subsection{Pipeline de données (Producer et Consumer)}

Le Producer supporte trois modes d'exécution~:

\begin{lstlisting}[language=Python,caption={Extrait du Producer tri-mode}]
# Mode 1: Historique (2020-2024)
def produce_historical(cities, start_date, end_date):
    for city in cities:
        data = fetch_open_meteo_historical(city, start_date, end_date)
        for record in data:
            producer.send('weather-historical', value=record)

# Mode 2: Temps reel (toutes les 15 min)
def produce_current(cities):
    for city in cities:
        data = fetch_owm_current(city)
        producer.send('weather-current', value=data)

# Mode 3: Previsions 7 jours
def produce_forecast(cities):
    for city in cities:
        data = fetch_open_meteo_forecast(city)
        producer.send('weather-forecast', value=data)
\end{lstlisting}

\begin{figure}[H]
\centering
\imgplaceholder{Figure~\ref{fig:producer_run} — Exécution du Producer}{Capture d'écran du terminal montrant le Producer envoyant des messages vers les topics Kafka}
\caption{Exécution du Producer — Envoi des données historiques}
\label{fig:producer_run}
\end{figure}

\subsection{API REST FastAPI}

\begin{figure}[H]
\centering
\imgplaceholder{Figure~\ref{fig:swagger} — Swagger UI}{Capture d'écran de la documentation Swagger de l'API FastAPI avec les endpoints /summary, /forecast, /auth}
\caption{Documentation Swagger de l'API FastAPI}
\label{fig:swagger}
\end{figure}

\subsection{Interface UI (Dashboard)}

\begin{figure}[H]
\centering
\imgplaceholder{Figure~\ref{fig:dashboard_final} — Dashboard final Sprint 1}{Capture d'écran du dashboard avec cartes KPI, carte Leaflet et graphiques Plotly.js en thème sombre}
\caption{Dashboard Météo Tunisie — Sprint 1}
\label{fig:dashboard_final}
\end{figure}

\section{Tests et Validation du Sprint 1}

\begin{table}[H]
\centering
\caption{Résultats des tests — Sprint 1}
\label{tab:tests_sprint1}
\begin{tabularx}{\textwidth}{clXc}
\toprule
\textbf{ID} & \textbf{Type} & \textbf{Description} & \textbf{Résultat} \\
\midrule
TC-01 & Unitaire     & Validation du token JWT                        & \ding{51} \\
TC-02 & Intégration  & Persistance des données après envoi Kafka       & \ding{51} \\
TC-03 & Fonctionnel  & Affichage des KPIs sur le dashboard              & \ding{51} \\
TC-04 & Performance  & Ingestion de 221 villes en $<$30s               & \ding{51} \\
TC-05 & Sécurité     & Connexion SSL Kafka avec certificats Aiven       & \ding{51} \\
\bottomrule
\end{tabularx}
\end{table}

% CHAPITRE 5 — SPRINT 2
\chapter{Sprint 2~: Analyses Avancées, Alertes et Déploiement}

\section{Introduction}
Le second sprint transforme les données brutes en informations stratégiques via des analyses statistiques et un système d'alertes automatisé.

\section{Objectif et Sprint Backlog}
\textbf{Objectif~:} Développer les modules d'analyses (corrélations, moyennes), le système de détection de pics, et déployer sur Astronomer Cloud.

\begin{table}[H]\centering\caption{Sprint Backlog — Sprint 2}\label{tab:sprint2_backlog}
\begin{tabularx}{\textwidth}{clX}\toprule
\textbf{US} & \textbf{Catégorie} & \textbf{Description} \\\midrule
US7 & Analyses & DAGs Airflow pour moyennes et corrélations \\
US8 & Alertes & Détection de pics thermiques (z-score) \\
US9 & Livraison & Documentation, monitoring, README \\\bottomrule
\end{tabularx}\end{table}

\section{Analyse et Spécification}

\subsection{Cas d'utilisation}
\begin{figure}[H]\centering
\imgplaceholder{Figure — Cas d'utilisation Sprint 2}{Diagramme UML~: Analyste et Système avec cas Visualiser Analyses, Détecter Pics, Recevoir Alertes}
\caption{Diagramme de cas d'utilisation — Sprint 2}\label{fig:uc_sprint2}\end{figure}

\begin{table}[H]\centering\caption{CU — Visualiser analyses statistiques (US7)}\label{tab:uc_stats}
\begin{tabularx}{\textwidth}{lX}\toprule
\textbf{Élément} & \textbf{Description} \\\midrule
Acteur & Analyste / Utilisateur \\
Scénario & 1. Sélection période \quad 2. Calcul corrélations Pearson \quad 3. Affichage Heatmap \\\bottomrule
\end{tabularx}\end{table}

\begin{table}[H]\centering\caption{CU — Détecter un pic de chaleur (US8)}\label{tab:uc_peak}
\begin{tabularx}{\textwidth}{lX}\toprule
\textbf{Élément} & \textbf{Description} \\\midrule
Acteur & Système \\
Scénario & 1. DAG peak\_detection s'exécute \quad 2. Calcul $z = (T-\mu)/\sigma$ \quad 3. Si $z>1.5$ insertion alerte \quad 4. Sévérité~: modéré/élevé/extrême \\\bottomrule
\end{tabularx}\end{table}

\section{Conception}

\subsection{Diagramme de classes}
\begin{figure}[H]\centering
\imgplaceholder{Diagramme de classes}{Classes~: WeatherData, MonthlyAverage, Correlation, TemperaturePeak, AnnualStats, DataQualityLog, PipelineRun}
\caption{Diagramme de classes final}\label{fig:class_diagram}\end{figure}

\subsection{Diagramme de séquence — Alerte}
\begin{figure}[H]\centering
\imgplaceholder{Séquence alerte}{Airflow $\rightarrow$ peak\_detection $\rightarrow$ Supabase $\rightarrow$ z-score $\rightarrow$ INSERT peaks $\rightarrow$ Frontend}
\caption{Séquence — Détection de pic thermique}\label{fig:seq_alert}\end{figure}

\subsection{Diagramme d'activité — Pipeline}
\begin{figure}[H]\centering
\imgplaceholder{Activité pipeline}{Check Freshness $\rightarrow$ Quality $\rightarrow$ [Monthly || Peaks] $\rightarrow$ Correlations $\rightarrow$ Annual $\rightarrow$ Complete}
\caption{Activité — master\_weather\_pipeline}\label{fig:activity_pipeline}\end{figure}

\section{Implémentation}

\subsection{DAGs Airflow}
\begin{table}[H]\centering\caption{DAGs Airflow}\label{tab:dags}
\begin{tabularx}{\textwidth}{lXl}\toprule
\textbf{DAG ID} & \textbf{Fonction} & \textbf{Schedule} \\\midrule
monthly\_analysis & Moyennes mensuelles & 07h quotidien \\
peak\_detection & Pics thermiques z-score & 07h quotidien \\
correlation\_analysis & Pearson saisonnières & Lundi 08h \\
annual\_stats & Stats annuelles, tendances & 1er janvier \\
data\_quality\_monitor & Contrôle qualité & 10h quotidien \\
master\_weather\_pipeline & Orchestration & 06h quotidien \\\bottomrule
\end{tabularx}\end{table}

\begin{lstlisting}[language=Python,caption={Architecture modulaire des DAGs}]
# dags/src/analysis/peaks.py
def detect_peaks_fn():
    df = pd.read_sql("SELECT * FROM weather_historical", engine)
    df['z_score'] = (df['temperature'] - df['temperature'].mean()) / df['temperature'].std()
    peaks = df[df['z_score'] > 1.5]
    return {'total': len(peaks)}
\end{lstlisting}

\begin{figure}[H]\centering
\imgplaceholder{Interface Airflow}{Capture des 6 DAGs avec statut et schedule}
\caption{Interface Airflow — 6 DAGs}\label{fig:airflow_ui}\end{figure}

\begin{figure}[H]\centering
\imgplaceholder{Graphe Master DAG}{Graphe du master\_weather\_pipeline}
\caption{Graphe master\_weather\_pipeline}\label{fig:dag_graph}\end{figure}

\subsection{Déploiement Astronomer Cloud}
\begin{lstlisting}[caption={Dockerfile Astronomer}]
FROM astrocrpublic.azurecr.io/runtime:3.2-3
\end{lstlisting}

\begin{figure}[H]\centering
\imgplaceholder{Astronomer Cloud}{Capture du déploiement réussi}
\caption{Déploiement Astronomer Cloud}\label{fig:astro_deploy}\end{figure}

\subsection{Interfaces Plotly.js}
\begin{figure}[H]\centering
\imgplaceholder{Heatmap corrélations}{Matrice de corrélation Plotly.js}
\caption{Matrice de corrélation saisonnière}\label{fig:heatmap}\end{figure}

\begin{figure}[H]\centering
\imgplaceholder{Pics thermiques}{Graphique des pics par sévérité}
\caption{Pics thermiques détectés}\label{fig:peaks_chart}\end{figure}

\section{Tests Sprint 2}
\begin{table}[H]\centering\caption{Tests Sprint 2}\label{tab:tests_sprint2}
\begin{tabularx}{\textwidth}{clXc}\toprule
\textbf{ID} & \textbf{Type} & \textbf{Description} & \textbf{OK} \\\midrule
TC-06 & Unitaire & Corrélations Pearson & \ding{51} \\
TC-07 & Unitaire & Pics z-score $>1.5$ & \ding{51} \\
TC-08 & Intégration & Master DAG complet & \ding{51} \\
TC-09 & Fonctionnel & Heatmap interactive & \ding{51} \\
TC-10 & Performance & DAG $<60$s / 221 villes & \ding{51} \\\bottomrule
\end{tabularx}\end{table}

% CHAPITRE 6 — TESTS GLOBAUX
\chapter{Tests Globaux et Validation Finale}

\section{Stratégie globale de test}
La stratégie de test couvre quatre niveaux~: tests unitaires, tests d'intégration, tests de qualité des données et tests de performance. L'objectif est de garantir la fiabilité du pipeline de bout en bout.

\section{Tests unitaires consolidés (pytest)}

\begin{lstlisting}[language=bash,caption={Exécution des tests unitaires}]
$ pytest tests/ -v --cov=src --cov-report=term-missing
======================== 125 passed in 12.4s ========================
\end{lstlisting}

\begin{table}[H]\centering\caption{Couverture des tests unitaires}\label{tab:coverage}
\begin{tabularx}{\textwidth}{Xcccc}\toprule
\textbf{Module} & \textbf{Tests} & \textbf{Passés} & \textbf{Couverture} \\\midrule
\texttt{src/utils/db.py} & 12 & 12 & 89\% \\
\texttt{src/utils/cities.py} & 18 & 18 & 95\% \\
\texttt{src/utils/validator.py} & 15 & 15 & 92\% \\
\texttt{src/utils/kafka\_config.py} & 10 & 10 & 88\% \\
\texttt{src/producer.py} & 35 & 35 & 85\% \\
\texttt{src/consumer.py} & 20 & 20 & 82\% \\
\texttt{src/utils/weather\_codes.py} & 8 & 8 & 100\% \\
\midrule
\textbf{Total} & \textbf{125} & \textbf{125} & \textbf{87\%} \\\bottomrule
\end{tabularx}\end{table}

\begin{figure}[H]\centering
\imgplaceholder{Résultats pytest}{Capture d'écran du terminal montrant les 125 tests passés avec couverture}
\caption{Résultats des tests unitaires — pytest}\label{fig:pytest}\end{figure}

\section{Tests d'intégration (pipeline E2E)}
Le test d'intégration valide le flux complet~: Kafka $\rightarrow$ DB $\rightarrow$ API $\rightarrow$ Frontend.

\begin{table}[H]\centering\caption{Tests d'intégration E2E}\label{tab:e2e}
\begin{tabularx}{\textwidth}{clXc}\toprule
\textbf{ID} & \textbf{Flux} & \textbf{Description} & \textbf{OK} \\\midrule
E2E-01 & Historique & Producer envoie 221 villes $\rightarrow$ Consumer persiste $\rightarrow$ API retourne & \ding{51} \\
E2E-02 & Temps réel & OWM API $\rightarrow$ Kafka $\rightarrow$ DB $\rightarrow$ Dashboard refresh & \ding{51} \\
E2E-03 & Prévisions & Open-Meteo $\rightarrow$ Kafka $\rightarrow$ DB $\rightarrow$ /forecast endpoint & \ding{51} \\
E2E-04 & Analyses & Airflow DAG $\rightarrow$ Supabase $\rightarrow$ API $\rightarrow$ Plotly.js & \ding{51} \\\bottomrule
\end{tabularx}\end{table}

\section{Tests de qualité des données}
Le DAG \texttt{data\_quality\_monitor} vérifie quotidiennement~:
\begin{itemize}
  \item \textbf{Complétude}~: absence de valeurs NULL sur les champs critiques
  \item \textbf{Validité}~: température dans $[-5°C, 55°C]$
  \item \textbf{Unicité}~: absence de doublons (date, city)
  \item \textbf{Fraîcheur}~: dernière ingestion $< 48$h
\end{itemize}

\begin{figure}[H]\centering
\imgplaceholder{Quality scores}{Capture montrant les scores de qualité par couche (historical, current, forecast)}
\caption{Scores de qualité des données}\label{fig:quality}\end{figure}

\section{Tests de performance}

\begin{table}[H]\centering\caption{Métriques de performance}\label{tab:perf}
\begin{tabularx}{\textwidth}{Xlc}\toprule
\textbf{Métrique} & \textbf{Valeur mesurée} & \textbf{Objectif} \\\midrule
Ingestion historique (221 villes $\times$ 5 ans) & 25s & $<$ 60s \\
Ingestion temps réel (221 villes) & 8s & $<$ 30s \\
Exécution DAG monthly\_analysis & 12s & $<$ 60s \\
Exécution DAG peak\_detection & 18s & $<$ 60s \\
Latence API /summary & 120ms & $<$ 500ms \\
Chargement Dashboard & 1.2s & $<$ 3s \\\bottomrule
\end{tabularx}\end{table}

\section{Résultats et métriques — Bilan de vélocité}

\begin{table}[H]\centering\caption{Vélocité de l'équipe}\label{tab:velocity}
\begin{tabularx}{\textwidth}{Xccc}\toprule
\textbf{Sprint} & \textbf{Points planifiés} & \textbf{Points livrés} & \textbf{Vélocité} \\\midrule
Sprint 1 & 50 & 50 & 100\% \\
Sprint 2 & 26 & 26 & 100\% \\\midrule
\textbf{Total} & \textbf{76} & \textbf{76} & \textbf{100\%} \\\bottomrule
\end{tabularx}\end{table}

\begin{figure}[H]\centering
\imgplaceholder{Burndown chart}{Graphique burndown montrant la progression des deux sprints}
\caption{Burndown Chart — Sprints 1 et 2}\label{fig:burndown}\end{figure}


% ─────────────────────────────────────────
%  CONCLUSION
% ─────────────────────────────────────────
\chapter*{Conclusion et Perspectives}
\addcontentsline{toc}{chapter}{Conclusion et Perspectives}

\section*{Bilan du projet}
Au terme de ce projet, nous avons conçu et déployé une plateforme Big Data complète dédiée à la météorologie tunisienne. Le système couvre l'intégralité du cycle de vie de la donnée~: collecte depuis les APIs Open-Meteo et OpenWeatherMap, ingestion en temps réel via Apache Kafka (Aiven), stockage persistant sur Supabase (PostgreSQL), orchestration analytique via Apache Airflow déployé sur Astronomer Cloud, et visualisation interactive à travers un dashboard responsive.

Les principaux résultats incluent~:
\begin{itemize}
  \item Couverture de \textbf{221 villes tunisiennes} avec un historique de 5 ans (2020--2024).
  \item \textbf{6 DAGs Airflow} opérationnels~: analyses mensuelles, annuelles, détection de pics, corrélations saisonnières, qualité des données et pipeline maître.
  \item \textbf{125 tests unitaires} validant l'ensemble des modules.
  \item Un dashboard interactif avec thème sombre/clair et cartes Leaflet.js.
\end{itemize}

\section*{Limites et difficultés rencontrées}
\begin{itemize}
  \item La gestion des connexions IPv6 sur le tier gratuit de Supabase a nécessité le recours au Session Pooler.
  \item La synchronisation des certificats SSL Kafka entre l'environnement local et le conteneur Docker Astronomer.
  \item Le volume de données (221 villes $\times$ 5 ans $\times$ 4 variables) impose des optimisations SQL pour les analyses.
\end{itemize}

\section*{Perspectives d'amélioration}
\begin{itemize}
  \item Intégration de modèles de Machine Learning (LSTM, Prophet) pour la prévision météo.
  \item Ajout d'un système de notifications push (Firebase) pour les alertes météo critiques.
  \item Extension aux données satellite (Copernicus) pour une couverture spatiale plus fine.
  \item Migration vers TimescaleDB pour optimiser les requêtes temporelles.
\end{itemize}

% ─────────────────────────────────────────
%  BIBLIOGRAPHIE
% ─────────────────────────────────────────
\chapter*{Bibliographie}
\addcontentsline{toc}{chapter}{Bibliographie}

\begin{enumerate}[label={[\arabic*]}]
  \item Apache Software Foundation, \textit{Apache Kafka Documentation}, \url{https://kafka.apache.org/documentation/}, consulté en 2025.
  \item Apache Software Foundation, \textit{Apache Airflow Documentation}, \url{https://airflow.apache.org/docs/}, consulté en 2025.
  \item Supabase Inc., \textit{Supabase Documentation}, \url{https://supabase.com/docs}, consulté en 2025.
  \item Aiven Ltd., \textit{Aiven for Apache Kafka}, \url{https://aiven.io/kafka}, consulté en 2025.
  \item Open-Meteo, \textit{Free Weather API}, \url{https://open-meteo.com/}, consulté en 2025.
  \item Astronomer Inc., \textit{Astronomer Documentation}, \url{https://docs.astronomer.io/}, consulté en 2025.
  \item S.~Ramírez, \textit{FastAPI Documentation}, \url{https://fastapi.tiangolo.com/}, consulté en 2025.
  \item W.~McKinney, \textit{Python for Data Analysis}, O'Reilly Media, 3\textsuperscript{e} édition, 2022.
  \item K.~Schwaber et J.~Sutherland, \textit{The Scrum Guide}, 2020.
  \item M.~Kleppmann, \textit{Designing Data-Intensive Applications}, O'Reilly Media, 2017.
\end{enumerate}

% ─────────────────────────────────────────
%  ANNEXES
% ─────────────────────────────────────────
\appendix
\chapter{Schéma SQL Complet}
\label{app:sql}

\begin{lstlisting}[language=SQL,caption={Extrait du schéma de la base de données Supabase}]
-- Table principale des mesures historiques
CREATE TABLE weather_historical (
    id            SERIAL PRIMARY KEY,
    date          DATE NOT NULL,
    city          VARCHAR(100) NOT NULL,
    governorate   VARCHAR(100),
    region        VARCHAR(50),
    temperature   REAL,
    humidity      REAL,
    precipitation REAL,
    wind_speed    REAL,
    ingested_at   TIMESTAMPTZ DEFAULT NOW(),
    UNIQUE(date, city)
);

-- Table des moyennes mensuelles (calculee par Airflow)
CREATE TABLE monthly_averages (
    id            SERIAL PRIMARY KEY,
    year          INTEGER NOT NULL,
    month         INTEGER NOT NULL,
    city          VARCHAR(100) NOT NULL,
    governorate   VARCHAR(100),
    region        VARCHAR(50),
    avg_temp      REAL,
    max_temp      REAL,
    min_temp      REAL,
    avg_humidity  REAL,
    avg_precip    REAL,
    avg_wind      REAL,
    record_count  INTEGER,
    computed_at   TIMESTAMPTZ DEFAULT NOW(),
    UNIQUE(year, month, city)
);

-- Table des correlations saisonnieres
CREATE TABLE correlations (
    id              SERIAL PRIMARY KEY,
    variable_a      VARCHAR(50) NOT NULL,
    variable_b      VARCHAR(50) NOT NULL,
    season          VARCHAR(20) NOT NULL,
    pearson_r       REAL,
    p_value         REAL,
    interpretation  VARCHAR(50),
    is_significant  BOOLEAN,
    computed_at     TIMESTAMPTZ DEFAULT NOW(),
    UNIQUE(variable_a, variable_b, season)
);
\end{lstlisting}

\chapter{Configuration Kafka et Variables d'Environnement}
\label{app:config}

\begin{lstlisting}[caption={Fichier .env.example}]
# Aiven Kafka (SSL)
KAFKA_BOOTSTRAP=kafka-projetfedere-....aivencloud.com:21849
KAFKA_SECURITY_PROTOCOL=SSL
KAFKA_CA_FILE=ca.pem
KAFKA_CERT_FILE=service.cert
KAFKA_KEY_FILE=service.key

# Supabase PostgreSQL
DB_HOST=aws-1-eu-central-1.pooler.supabase.com
DB_PORT=5432
DB_NAME=postgres
DB_USER=postgres.your_project_ref
DB_PASSWORD=your_password

# Weather APIs
TIMEZONE=Africa/Tunis
START_DATE=2020-01-01
END_DATE=2024-12-31
OWM_API_KEY=your_api_key
\end{lstlisting}

\chapter{Guide de Déploiement}
\label{app:deploy}

\section{Supabase}
\begin{enumerate}
  \item Créer un projet sur \url{https://supabase.com}
  \item Exécuter le schéma SQL (Annexe~\ref{app:sql}) dans l'éditeur SQL
  \item Récupérer le host du Session Pooler (Settings $\rightarrow$ Database)
\end{enumerate}

\section{Aiven Kafka}
\begin{enumerate}
  \item Créer un service Kafka sur \url{https://aiven.io}
  \item Télécharger les certificats SSL (ca.pem, service.cert, service.key)
  \item Créer les topics~: \texttt{weather-historical}, \texttt{weather-current}, \texttt{weather-forecast}
\end{enumerate}

\section{Astronomer Cloud (Airflow)}
\begin{enumerate}
  \item Installer Docker Desktop et le CLI Astro~: \texttt{brew install astro}
  \item Se connecter~: \texttt{astro login}
  \item Valider les DAGs~: \texttt{astro dev parse}
  \item Déployer~: \texttt{astro deploy}
  \item Configurer les variables d'environnement dans l'UI Astronomer
\end{enumerate}

\chapter{Manuel Utilisateur}
\label{app:user}

\section{Accès au Dashboard}
L'utilisateur accède à la plateforme via un navigateur web. La page d'accueil présente un tableau de bord avec les indicateurs météo principaux.

\imgplaceholder{Figure — Dashboard principal}{Capture d'écran du dashboard avec les KPIs météo, la carte Leaflet et les graphiques Plotly.js}

\section{Navigation}
\begin{itemize}
  \item \textbf{Accueil}~: Vue d'ensemble avec carte interactive
  \item \textbf{Temps réel}~: Données actuelles avec rafraîchissement automatique
  \item \textbf{Prévisions}~: Prévisions à 7 jours par ville
  \item \textbf{Analyses}~: Corrélations, tendances annuelles, pics thermiques
\end{itemize}

\end{document}
